\documentclass[10pt]{article}

\usepackage[utf8]{inputenc}

\usepackage[T1]{fontenc}

\usepackage{amsmath}

\usepackage{amsfonts}

\usepackage{amssymb}

\usepackage[version=4]{mhchem}

\usepackage{stmaryrd}

\usepackage{bbold}



\begin{document}

причем обычно предполагается, что



\[

\begin{aligned}

& F=\operatorname{grad} f, f: \mathbb{R}^{n} \longrightarrow \mathbb{R}, f \in C^{\mathbf{1}}, \\

& G=\operatorname{grad} g, g: \mathbb{R}^{n} \longrightarrow \mathbb{R}, g \in C^{2},

\end{aligned}

\]



a $H$ - ограниченная и периодическая по $t$ векторфункция. Представляет интерес получение достаточных условий существования периодич. решений таких систем.



Лum.: [1] С т р е т т Дж. В. (лор д Р ә е й), Теория звука, пер. с англ., 2 изд., т. 1, М.- ЈІ., 1955, [2] ч е за р и Л., Асимптотическое поведение и устойчивость решений обыкновенных дифференциальных уравнений, пер. с англ., М., 1964. См. РЯД, бе с к о неч н а я с у м м а,- последовательность элементов (наз. ч л е н а м и д а н н о г о р я д а) нек-рого линейного топологич. простракства и определенное бесконечное множество их конечных сумм (наз. ч а с т и ч ны м и с у м м а ми р я д а), для к-рых определено понятие предела. Простейшими примерами Р. являются следующие.



Однократные числовые ряды. Пара последовательностей комплексных чисел $\left\{a_{n}\right\}$ и $\left\{s_{n}\right\}$ таких, что



\[

s_{n}=a_{1}+a_{2}+\ldots+a_{n}, n=1,2, \ldots,

\]



наз. ч и с ло в ы (о днок р а т ны м) р я до м и обозначается так:



\[

a_{1}+a_{2}+\ldots+a_{n}+\ldots, \sum_{n=1}^{\infty} a_{n}

\]



или



\[

\sum a_{n}

\]



Элементы последовательности $\left\{a_{n}\right\}$ наз. ч ле н а м и р я д а, а элементы последовательности $\left\{s_{n}\right\}$ - его ч а с т и ч ны м и с у м м а м и, причем $a_{n}$ наз. $n-\mathbf{M}$ членом P. (2), а $s_{n}$ - его частичной суммой порядка $n$. Р. (2) однозначно определяется каждой из двух последовательностей $\left\{a_{n}\right\}$ и $\left\{s_{n}\right\}:$ члены последовательности $\left\{s_{n}\right\}$ получаются из членов последовательности $\left\{a_{n}\right\}$ по формуле (1), а последовательность $\left\{a_{n}\right\}$ восстанавливается по последовательности $\left\{s_{n}\right\}$ согласно формулам



\[

a_{1}=s_{1}, a_{n+1}=s_{n+1}-s_{n}, \quad n=1,2, \ldots

\]



В этом смысле изучение Р. равносильно изучению последовательностей: для каждого утверждения о Р. можно сформулировать равносильное ему утверждение о последовательностях.



Р. (2) наз. с х о д я п и м с я, если последовательность его частичных сумм $\left\{s_{n}\right\}$ имеет конечный предел



\[

s=\lim _{n \rightarrow \infty} s_{n}

\]



к-рый наз. с у м м о й Р. (2), и пишут



\[

s=\sum a_{n}

\]



Таким образом, обозначение (2) применяется как для самого Р., так и для его суммы. Если последовательность частичных сумм Р. (2) не имеет конечного предела, то он наз. $\mathrm{p}$ а с х о д я Щ и с я.



Примером сходящегося Р. является сумма членов бесконечной геометрич. прогрессии



\[

\sum_{n=1}^{\infty} q^{n}

\]



шри условии, что $|q|<1$. В этом случае ее сумма равна $\frac{1}{1-q}$, то есть $\sum_{n=1}^{\infty} q^{n}=\frac{1}{1-q}$. Если же $|q| \geqslant 1$, то Р. (3) дает пример расходящегося Р.



Если P (2) сходится, то последовательность его членов стренится к нулю:



\[

\lim _{n \rightarrow \infty} a_{n}=0

\]



яд



Обратное утверждение неверно: последовательность членов $\left\{\frac{1}{n}\right\}$ гармонического ряда



\[

1+\frac{1}{2}+\frac{1}{3}+\ldots+\frac{1}{n}+\ldots

\]



стремится к нулю, однако этот Р. расходится. Ряд $\sum_{k=1}^{\infty} a_{n+k}$ наз. о с т а т к о м порядка $n$ Р. (2). Если P. сходится, то каждый его остаток сходится. Если нек-рый остаток Р. сходится, то и сам Р. сходится. Если остаток порядка $n$ P. (2) сходится и его сумма равна $r_{n}$, то есть $r_{n}=\sum_{k=1}^{\infty} a_{n+k}$, то



\[

s=\sum_{n=1}^{\infty} a_{n}=s_{n}+r_{n} .

\]



Если Р. (2) и ряд



$\sum b_{n}$



сходятся, то сходится и ряд



\[

\sum\left(a_{n}+b_{n}\right)

\]



наз. с у м м о й рядов (2) и (3), причем его сумма равна сумме этих $P$.



Если Р. (2) сходится и $\lambda$-комплексное число, то ряд $\sum \lambda a_{n}$, наз. п р о и з в е д е н и е м Р. (2) на число $\lambda$, также сходится и $\sum \lambda a_{n}=\lambda \sum a_{n}$.



Условие сходимости Р., не использующее понятие его суммы, дает Коши критерий сходимости Р.



Если все члены P. (2) являются действительными числами: $a_{n} \in \mathbb{R}$, то Р. (2) наз. действительным. Важную роль в теории Р. играют действительные Р. с неотрицательными членами:



\[

\sum a_{n}, a_{n} \geqslant 0, n=1,2, \ldots .

\]



Для того чтобы Р. (5) сходился, необходимо и достаточно, чтобы последовательность его частичных сумм была ограничена сверху. Если же он расходится, то его частичные суммы стремятся к бесконечности:



\[

\lim _{n \rightarrow \infty} s_{n}=+\infty

\]



поэтому в этом случае пишут



\[

\sum a_{n}=+\infty

\]



Для Р. с неотрицательными членами существует много различных признаков сходимости. Основными являются следующие.



Признак с равнения. Если для Р. (5) и



\[

\sum b_{n}, b_{n} \geqslant 0, \quad n=1,2, \ldots,

\]



с неотрицательными членами существует такая постоянная $c>0$, что $0 \leqslant a_{n} \leqslant c b_{n}$, то из сходимости $\mathrm{P}$. (6) следует сходимость Р. (5), а из расходимости Р. (5) - расходимость Р. (6).



При применении признака сравнения для исследования сходимости заданного Р. с неотрицательными членами часто оказывается целесообразным выделить главную часть его $n$-го члена относительно $\frac{1}{n}$ при $n \rightarrow \infty$ в виде $\frac{a}{n^{\alpha}} \quad(a$ - нек-рая постоянная), а в качестве Р., с к-рым сравнивается данный Р., взять ряд



\[

\sum_{n=1}^{\infty} \frac{1}{n^{\alpha}}, \alpha \in \mathbb{R}

\]



сходящийся при $\alpha>1$ п расходящийся при $\alpha \leqslant 1$.





\end{document}